\subsection{APPROXIMATION OF CONTINUOUS FUNCTIONS BY FUNCTIONS BELONGING TO A LATTICE}

In this section we shall study the set \(\mathcal{C} = \mathcal{C}(\mathbf{X};\mathbf{R})\) of continuous real-valued functions (*) defined on a compact space \(\mathbf{X}\) , and we shall always suppose that \(\mathcal{C}\) is endowed with the topology of uniform convergence. From § 3, no. 2 we know that this topology is defined by the norm

\[
| | f | | = \sup _ {x \in \mathbf {X}} | f (x) |
\]

and that this norm is compatible with the R-algebra structure of C. With this norm and this algebra structure, C is a complete normed algebra over R (§ 1, no. 6, Theorem 2, Corollary 1).

If H is a subset of C, we shall say that a continuous real-valued function f on X can be uniformly approximated by functions of H if f lies in the closure of H in the space C, i.e.~if, for each \(\varepsilon > 0\) , there exists a function \(g \in H\) such that \(|f(x) - g(x)| \leqslant \varepsilon\) for all \(x \in X\) . To say that every continuous real-valued function on X can be uniformly approximated by functions of H therefore means that H is dense in C.

On the set \(\mathcal{C}\) , the relation \(f \leqslant g\) {[}which means that \(f(x) \leqslant g(x)\) for all \(x \in \mathbf{X}\) {]} is an order relation, with respect to which \(\mathcal{C}\) is a lattice. Clearly we have \(|||u| - |v|| \leqslant ||u - v||\) , and therefore \(u \to |u|\) is a uniformly continuous mapping of \(\mathcal{C}\) into itself. It follows that

\[
(u, v) \rightarrow \sup (u, v) = \frac {1}{2} (u + v + | u - v |)
\]

and

\[
(u, v) \rightarrow \inf (u, v) = \frac {1}{2} (u + v - | u - v |)
\]

are uniformly continuous on \(\mathcal{C} \times \mathcal{C}\) .

PROPOSITION 1. Let X be a compact space and let H be a set of continuous real-valued functions defined on X. Let f be a continuous real-valued function on X such that for each \(x \in X\) there exists a function \(u_x \in H\) such that \(u_x(x) > f(x)\) {[}resp. \(u_x(x) < f(x)\) {]}. Then there exists a finite number of functions \(u_{x_i} = f_i \in H\) ( \(1 \leqslant i \leqslant n\) ) such that, if \(v = \sup(f_1, f_2, \ldots, f_n)\) {[}resp. \(w = \inf(f_1, f_2, \ldots, f_n)\) {]}, we have \(v(x) > f(x)\) {[}resp. \(w(x) < f(x)\) {]} for all \(x \in X\) .

For each \(x \in \mathbf{X}\) , let \(\mathbf{G}_x\) be the open set consisting of all \(z \in \mathbf{X}\) such that \(u_x(z) > f(z)\) {[}resp. \(u_x(z) < f(z)\) {]}. Since \(x \in \mathbf{G}_x\) by hypothesis, \(\mathbf{X}\) is the union of the sets \(\mathbf{G}_x\) as \(x\) runs through \(\mathbf{X}\) . Since \(\mathbf{X}\) is compact there exists a finite number of points \(x_i (1 \leqslant i \leqslant n)\) such that the \(\mathbf{G}_{x_i}\) cover \(\mathbf{X}\) , and it is clear that the functions \(f_i = u_{x_i}\) satisfy the conditions of the proposition.

THEOREM 1 (Dini). Let X be a compact space, and let H be a set of continuous real-valued functions on X which is directed with respect to the relation \(\leqslant (\text{resp.} \geqslant)\) . If the upper (resp. lower) envelope f of H is finite and continuous on X, then f can be uniformly approximated by functions belonging to H (or, equivalently, the section filter of H converges uniformly to f in X).

Given any \(\varepsilon > 0\) , for each \(x \in \mathbf{X}\) there exists a function \(u_x \in \mathbf{H}\) such that \(u_x(x) > f(x) - \varepsilon\) . By Proposition 1 and the fact that H is directed with respect to the relation \(\leqslant\) , there exists \(g \in \mathbf{H}\) such that \(g(x) > f(x) - \varepsilon\) for all \(x \in \mathbf{X}\) ; on the other hand, we have \(g(x) \leqslant f(x)\) by definition, and therefore the theorem is proved.

COROLLARY. Let \((u_{n})\) be an increasing (resp. decreasing) sequence of continuous real-valued functions on X. If the upper (resp. lower) envelope f of the sequence \((u_{n})\) is finite and continuous on X, then the sequence \((u_{n})\) converges uniformly to f in X.

It is clear that the conclusion of Theorem 1 is no longer necessarily valid if X is no longer assumed to be compact, as is shown by the example of the decreasing sequence of functions \(x/(n+x)\) in \(R_{+}\) .

PROPOSITION 2. Let X be a compact space, and let H be a set of continuous real-valued functions on X such that, given any two functions \(u \in H\) , \(v \in H\) , the functions \(\sup(u, v)\) and \(\inf(u, v)\) are in H. Then a continuous real-valued function \(f\) on X can be uniformly approximated by functions belonging to H if and only if, for each real number \(\varepsilon > 0\) and each pair \(x, y\) of points of X, there is a function \(u_{x,y} \in H\) such that \(|f(x) - u_{x,y}(x)| < \varepsilon\) and \(|f(y) - u_{x,y}(y)| < \varepsilon\) .

The condition is clearly necessary; let us show that it is sufficient. For each \(\varepsilon > 0\) , we shall show that there is a function \(g \in \mathbf{H}\) such that \(|f(z) - g(z)| < \varepsilon\) for all \(z \in \mathbf{X}\) . Let \(x\) be any point of \(\mathbf{X}\) , and let \(\mathbf{H}_x\) be the set of all functions \(u \in \mathbf{H}\) such that \(u(x) < f(x) + \varepsilon\) . By hypothesis, for each \(y \in \mathbf{X}\) , the function \(u_{x,y}\) belongs to \(\mathbf{H}_x\) and we have \(u_{x,y}(y) > f(y) - \varepsilon\) . Hence, by Proposition 1, there is a finite number of functions of \(\mathbf{H}_x\) whose upper envelope \(v_x\) is such that \(v_x(z) > f(z) - \varepsilon\) for all \(z \in \mathbf{X}\) ; on the other hand, we have \(v_x(x) < f(x) + \varepsilon\) by the definition of \(\mathbf{H}_x\) ; finally, \(v_x \in \mathbf{H}\) by hypothesis. Proposition 1 therefore shows that there exists a finite number of functions \(v_{x_i}\) whose lower envelope \(g\) is such that \(g(z) < f(z) + \varepsilon\) for all \(z \in \mathbf{X}\) ; but since we have \(v_{x_i}(z) > f(z) - \varepsilon\) for all \(z \in \mathbf{X}\) and for every index \(i\) , we have also \(g(z) > f(z) - \varepsilon\) for all \(z \in \mathbf{X}\) . Since \(g \in \mathbf{H}\) by hypothesis, the proof is complete.

\pandocbounded{\includegraphics[keepaspectratio]{images/part002_45bf5402.jpg}}

Remark. When the set H satisfies the conditions of Proposition 2, it is a lattice with respect to the ordering \(f \leqslant g\) . But it should be remarked that a subset H of C can be a lattice with respect to this ordering without it being necessarily the case that the least upper bound (resp. greatest lower bound) in H of two functions u, v of H is the same as their least upper bound (resp. greatest lower bound) in C. * An example is provided by the convex mappings of a compact interval of R into R.*

COROLLARY. Suppose that H is such that, whenever \(u \in H\) and \(v \in H\) , we have sup \((u, v) \in H\) and inf \((u, v) \in H\) ; and is such that, given any two distinct points x, y of X and any two real numbers \(\alpha, \beta\) , there is a function \(g \in H\) such that \(g(x) = \alpha\) and \(g(y) = \beta\) . Then every continuous real-valued function on X can be uniformly approximated by functions belonging to H.

DEFINITION 1. If X is any set, a set H of mappings of X into a set Y is said to separate the elements of a subset A of X (or to be a separating set for the elements of A) if, given any two distinct elements x, y of A, there is a function \(f \in H\) such that \(f(x) \neq f(y)\) .

For example, if X is a completely regular space (Chapter IX, § 1, no. 5) then the set of all continuous mappings of X into {[}0, 1{]} separates the points of X.

THEOREM 2 (Stone). Let X be a compact space, and let H be a vector subspace of C(X; R) such that 1) the constant functions belong to H; 2) if \(u \in H\) , then \(|u| \in H\) ; 3) H separates the points of X. Then every continuous real-valued function on X can be uniformly approximated by functions of H.

It is enough to show that H satisfies the conditions of the Corollary to Proposition 2. By hypothesis, if \(u \in H\) and \(v \in H\) , we have

\[
\sup (u, v) = \frac {1}{2} (u + v + | u - v |) \in H
\]

and

\[
\inf (u, v) = \frac {1}{2} (u + v - | u - v |) \in H.
\]

On the other hand, let \(x\) and \(y\) be any two distinct points of \(X\) , and let \(\alpha, \beta\) be any two real numbers. By hypothesis, there is a function \(h \in H\) such that \(h(x) \neq h(y)\) : say \(h(x) = \gamma\) and \(h(y) = \delta\) . Since the constants belong to \(H\) , the function

\[
g (z) = \alpha + (\beta - \alpha) \frac {h (z) - \gamma}{\delta - \gamma}
\]

belongs to H and is such that \(g(x) = \alpha\) and \(g(y) = \beta\) .

